\documentclass[10pt,letterpaper]{article}
\usepackage[margin=0.68in]{geometry}
\usepackage{microtype}
\usepackage{amsmath,amssymb}
\usepackage{booktabs}
\usepackage{tabularx}
\usepackage{graphicx}
\usepackage[colorlinks=true,linkcolor=blue,urlcolor=blue,citecolor=blue]{hyperref}

\setlength{\parindent}{0pt}
\setlength{\parskip}{5pt}
\sloppy

\title{\textbf{Deterministic Sub-100 $\mu$s Bare-Metal FPGA Spatial Voxel Screening and Dual-Target Structural Discovery of Allosteric TYK2 JH2 Inhibitors Selective Over JAK2}}
\author{\textbf{Joseph Iko}$^{1}$ \\
\small $^{1}$Oikonomia BioPharma / Oikonomia Hardware Architecture Lab, New Jersey, USA \\
\small \textbf{Subject Area:} Bioinformatics / Biophysics and Computational Biology}
\date{}

\begin{document}
\maketitle
\vspace{-1.5em}

\section*{Abstract}
Selective allosteric inhibition of the tyrosine kinase 2 pseudokinase domain (TYK2 JH2) over the catalytic Janus kinase 2 domain (JAK2 JH1) avoids dose-limiting hematologic toxicities while treating autoimmune pathology. Here, we report a dual-target, structure-based virtual screening and multi-seed consensus docking investigation comparing human TYK2 JH2 (PDB: 6NZP, 2.35~Å) against off-target catalytic JAK2 JH1 (PDB: 4IVA, 1.50~Å; 1.887~Å\ conserved binding-site C$\alpha$ RMSD). From 2,500 unique chemical structures curated across ChEMBL~34, known JH2 controls, and structure-guided oral-window analogs, medicinal chemistry alert triage (PAINS, Brenk, NIH), physicochemical filtering ($280 \le \text{MW} \le 490\text{ Da}$), and hinge-binding pharmacophore matching retained 385 qualified three-dimensional (3D) conformers. De novo SMILES redocking of the clinical allosteric inhibitor deucravacitinib (BMS-986165) into 6NZP achieved $\Delta G = -8.733\text{ kcal/mol}$ (4.412~Å\ element-aware heavy-atom RMSD), while active-versus-decoy classification across known JH2 binders and clinical drug decoys achieved $\text{ROC-AUC} = 0.867$. Multi-seed AutoDock Vina consensus docking (seeds 42, 137, 2026) identified scaffold-hopped ChEMBL~34 leads including \textbf{CHEMBL6914} ($\Delta G_{\text{TYK2}} = -10.72 \pm 0.04\text{ kcal/mol}$, $\Delta\Delta G_{\text{TYK2-JAK2}} = -1.11\text{ kcal/mol}$, $\text{LLE} = 3.05$, Val690 = 3.34~Å, Lys642 = 2.86~Å) and \textbf{CHEMBL266324} ($\Delta G_{\text{TYK2}} = -10.11 \pm 0.12\text{ kcal/mol}$, $\Delta\Delta G = -1.79\text{ kcal/mol}$, $\text{QED} = 0.654$), alongside the oral-window analog \textbf{OIK-TYK2-003} ($\text{MW} = 421.46\text{ Da}$, $\text{LLE} = 4.14$, $\Delta G_{\text{TYK2}} = -8.84 \pm 0.02\text{ kcal/mol}$). To enforce cycle-accurate spatial tensor verification without host DRAM overhead, 4,096-byte ($16\times16\times16$) molecular pharmacophore grids were streamed via direct userspace PCIe Gen3 x16 BAR0 memory-mapped I/O (MMIO) to a Xilinx UltraScale+ FPGA (\texttt{agfi-0b55f06c15a385187} on AWS EC2 \texttt{f2.6xlarge}), achieving 87.03~$\mu$s P50 (102.998~$\mu$s P99) C BAR0 streaming latency, 0.25~KB constant $O(1)$ on-chip state, and 100\% bit-exact register determinism at 9~W average (10~W max) $V_{\text{CCINT}}$ core power.

\section{Introduction}
Tyrosine kinase 2 (TYK2) mediates intracellular cytokine signaling central to immune-mediated inflammatory diseases (IMIDs) including psoriasis, systemic lupus erythematosus, and inflammatory bowel disease. Because catalytic ATP-binding sites across the Janus kinase (JAK1, JAK2, JAK3, TYK2) family share high sequence and structural conservation, orthosteric pan-JAK inhibitors frequently cross-react with JAK2 JH1, triggering dose-limiting cytopenias and hematologic toxicities [1]. In contrast, targeting the regulatory pseudokinase domain (TYK2 JH2) stabilizes an autoinhibitory conformation via hydrogen-bonding interactions with hinge residue Val690 and catalytic-loop Lys642 while sparing JAK2 JH1 [1].

Despite the clinical validation of deuterated allosteric inhibitor deucravacitinib (BMS-986165), discovering lower-molecular-weight ($\text{MW} \le 490\text{ Da}$), synthetically tractable chemotypes with favorable lipophilic ligand efficiency ($\text{LLE}$) and isoform selectivity over JAK2 remains an active challenge in computational medicinal chemistry. Furthermore, evaluating hundreds of 3D molecular pharmacophore grids across conventional host CPU/GPU servers incurs substantial DRAM allocation overhead and OS thread preemption at 150--350~W socket power. Here, we couple a dual-target structural biology workflow (TYK2 JH2 vs.\ JAK2 JH1) with OIK-PAS-V1, a direct userspace PCIe Gen3 x16 FPGA spatial voxel substrate operating at 9~W average core power.

\section{Results}
\subsection{Virtual Library Triage and Dual-Target Selectivity Screen (TYK2 JH2 vs.\ JAK2 JH1)}
From 2,500 unique chemical structures ingested across ChEMBL~34 [3], literature JH2 positive controls, and structure-guided oral-window analogs, automated medicinal chemistry triage enforcing RDKit \texttt{FilterCatalog} removal of PAINS, Brenk, and NIH structural alerts, oral Lipinski physicochemical boundaries ($280 \le \text{MW} \le 490\text{ Da}$, $1.0 \le \text{cLogP} \le 4.8$, $45.0 \le \text{TPSA} \le 125.0\text{ Å}^2$, $\text{HBD} \le 4$, $\text{HBA} \le 9$, $\text{RotB} \le 8$), and presence of an aromatic hinge-binding hydrogen-bond motif (\texttt{[n,N]\textasciitilde[c]\textasciitilde[n,N]}) retained 385 qualified 3D conformers. Spatial-voxel-ranked candidates across Track~1 (Scaffold-Hopped ChEMBL~34 Hits) and Track~2 (Structure-Guided Oral-Window Analogs) were counter-screened against both TYK2 JH2 (6NZP) and off-target JAK2 JH1 (4IVA) and audited across three independent random seeds (42, 137, 2026):

\begin{table}[h!]
\centering
\resizebox{\textwidth}{!}{%
\begin{tabular}{@{}llrrrrrrrrrrr@{}}
\toprule
\textbf{Track/Rank} & \textbf{Compound ID} & \textbf{MW} & \textbf{cLogP} & \textbf{LLE} & \textbf{QED} & \textbf{SA} & \textbf{Tan.} & \textbf{$\Delta G_{\text{TYK2}}$} & \textbf{$\Delta G_{\text{JAK2}}$} & \textbf{$\Delta\Delta G$} & \textbf{Val690/Lys642 (Å)} & \textbf{3-Seed Consensus} \\
\midrule
Track 1 \#1 & CHEMBL6914 & 487.43 & 4.79 & 3.05 & 0.402 & 2.96 & 0.103 & -10.70 & -9.59 & -1.11 & 3.34 / 2.86 & $-10.72 \pm 0.04$ \\
Track 1 \#2 & CHEMBL414010 & 453.88 & 4.60 & 3.22 & 0.438 & 3.42 & 0.096 & -10.67 & -9.72 & -0.95 & 3.19 / 2.87 & $-10.61 \pm 0.07$ \\
Track 1 \#3 & CHEMBL7029 & 487.43 & 4.79 & 2.83 & 0.402 & 2.94 & 0.112 & -10.39 & -9.59 & -0.80 & 3.44 / 4.63 & $-10.44 \pm 0.05$ \\
Track 1 \#4 & CHEMBL415478 & 475.42 & 4.32 & 3.23 & 0.407 & 2.77 & 0.105 & -10.30 & -9.29 & -1.01 & 3.22 / 3.43 & $-10.38 \pm 0.07$ \\
Track 1 \#5 & CHEMBL266324 & 421.55 & 3.89 & 3.65 & 0.654 & 2.68 & 0.198 & -10.28 & -8.49 & -1.79 & 3.42 / 4.66 & $-10.11 \pm 0.12$ \\
Track 2 \#1 & OIK-TYK2-003 & 421.46 & 2.34 & 4.14 & 0.535 & 2.92 & 0.722 & -8.84 & -8.41 & -0.43 & 4.04 / 3.11 & $-8.84 \pm 0.02$ \\
\bottomrule
\end{tabular}%
}
\caption{Dual-target AutoDock Vina consensus docking ($\text{kcal/mol}$) of top voxel-screened candidates against human TYK2 JH2 (PDB: 6NZP) and off-target JAK2 JH1 (PDB: 4IVA) across seeds 42, 137, and 2026. Tanimoto similarity is computed against deucravacitinib (Morgan radius 2, 2048-bit).}
\end{table}

\subsection{Structural Alignment and Crystallographic Control Validation}
High-resolution crystal structures of human TYK2 JH2 co-crystallized with deucravacitinib (PDB: 6NZP, 2.35~Å) and off-target human JAK2 JH1 (PDB: 4IVA, 1.50~Å) were aligned via global pairwise sequence alignment and C$\alpha$ superposition across 18 conserved binding-site C$\alpha$ residues ($\le 12.5$~Å\ from crystallographic ligand LB7), yielding a 1.887~Å\ binding-site C$\alpha$ RMSD and establishing the allosteric JH2 pocket centroid at $(14.16, -3.51, 28.73)$~Å.

\begin{table}[h!]
\centering
\footnotesize
\begin{tabularx}{\textwidth}{@{}llrrX@{}}
\toprule
\textbf{Validation Benchmark} & \textbf{Protocol / Target} & \textbf{$\Delta G$} & \textbf{RMSD} & \textbf{Key Structural Contact \& Enrichment} \\
\midrule
Deucravacitinib Control & De Novo SMILES $\rightarrow$ 6NZP & $-8.733\text{ kcal/mol}$ & $4.412\text{ Å}$ & Allosteric hinge Val690 + Lys642 pocket engagement \\
Active vs.\ Decoy Separation & 3 JH2 Binders vs.\ 5 Decoys & -- & -- & $\text{ROC-AUC} = 0.867$ rank discrimination \\
\bottomrule
\end{tabularx}
\caption{Control validation across crystallographic reference inhibitor deucravacitinib (BMS-986165) and multi-ligand active-versus-decoy enrichment (3 known JH2 binders vs.\ aspirin, caffeine, ibuprofen, atorvastatin, and metformin).}
\end{table}

\subsection{Direct Userspace PCIe Gen3 x16 Silicon Telemetry}
To eliminate Linux kernel driver overhead and OS thread preemption, 3D molecular pharmacophore grids were encoded as 4,096-byte ($16\times16\times16\times 8\text{-bit}$) spatial tensors (with atomic number and H-bond donor/acceptor flags) and streamed directly to physical Xilinx UltraScale+ silicon (\texttt{0000:34:00.0}, Amazon FPGA Image \texttt{agfi-0b55f06c15a385187}) via direct userspace PCIe BAR0 MMIO pinned to isolated host core CPU~2 (\texttt{taskset -c 2}).

\begin{table}[h!]
\centering
\footnotesize
\begin{tabularx}{\textwidth}{@{}llX@{}}
\toprule
\textbf{Hardware / Telemetry Metric} & \textbf{Measured Value} & \textbf{Execution Profile \& Architectural Significance} \\
\midrule
FPGA Silicon \& Image ID & \texttt{agfi-0b55f06c15a385187} & Xilinx UltraScale+ (\texttt{0000:34:00.0}) on AWS EC2 \texttt{f2.6xlarge} \\
Core Power Rail ($V_{\text{CCINT}}$) & 9~W Avg / 10~W Max & Verified via \texttt{fpga-describe-local-image -S 0 -R -M} ($\approx 25\times$--$40\times$ below CPU TDP) \\
On-Chip State Footprint & 0.25~KB ($O(1)$ constant) & Zero host DRAM allocation or page-table thrashing \\
C BAR0 Stream P50 / P99 ($385 \times 4$~KB) & 87.03~$\mu$s / 102.998~$\mu$s & Native \texttt{liboikonomia\_voxel\_dma.so} (\texttt{gcc -O3 -march=native}) \\
Host CPU / GPU Socket Baseline & 150--350~W TDP & Conventional host memory evaluation subject to OS preemption \\
Hardware Output Determinism & 100\% Bit-Exact & Verified via repeated BAR0 register readback (\texttt{Determinism=True}) \\
\bottomrule
\end{tabularx}
\caption{Empirical physical silicon telemetry captured on AWS EC2 \texttt{f2.6xlarge} (\texttt{Ubuntu 26.04 LTS}) with nanosecond \texttt{CLOCK\_MONOTONIC\_RAW} instrumentation.}
\end{table}

\section{Discussion \& Limitations}
\begin{enumerate}
    \item \textbf{Scaffold-Hopped Chemotypes from ChEMBL~34}: In Track~1, \textbf{CHEMBL6914} ($\text{MW} = 487.43\text{ Da}$, $\text{LLE} = 3.05$, Morgan Tanimoto $= 0.103$ to deucravacitinib) achieves $-10.72 \pm 0.04\text{ kcal/mol}$ consensus affinity with $-1.11\text{ kcal/mol}$ selectivity over JAK2 JH1 while forming dual polar contacts with hinge Val690 (3.34~Å) and Lys642 (2.86~Å). \textbf{CHEMBL266324} ($\text{MW} = 421.55\text{ Da}$, $\text{QED} = 0.654$, $\text{LLE} = 3.65$, $\text{SA} = 2.68$) achieves $-10.11 \pm 0.12\text{ kcal/mol}$ affinity and $-1.79\text{ kcal/mol}$ selectivity over JAK2 JH1, providing a synthetically tractable, drug-like starting scaffold.
    \item \textbf{Oral-Window TYK2 JH2 Selectivity}: In Track~2, \textbf{OIK-TYK2-003} ($\text{MW} = 421.46\text{ Da}$, $\text{cLogP} = 2.34$, $\text{LLE} = 4.14$) achieves $-8.84 \pm 0.02\text{ kcal/mol}$ consensus affinity with favorable lipophilic efficiency and a 3.11~Å\ Lys642 polar anchor while sparing JAK2 JH1 ($\Delta\Delta G = -0.43\text{ kcal/mol}$).
    \item \textbf{Architectural Trade-Offs \& Low-Power Determinism}: Streaming 1,024 32-bit scalar MMIO writes per 4,096-byte tensor over PCIe BAR0 measures control-plane register ingress (87.03~$\mu$s P50, 102.998~$\mu$s P99). Rather than raw single-pose scalar throughput, the primary hardware contribution of OIK-PAS-V1 is 9~W average (10~W max) $V_{\text{CCINT}}$ steady-state power ($\approx 25\times$--$40\times$ lower than server CPU/GPU TDP), a constant $O(1)$ 0.25~KB on-chip state footprint, and 100\% bit-exact cycle determinism free of OS thread preemption.
\end{enumerate}

\textbf{Limitations \& Experimental Scope:} AutoDock Vina [4, 5] treats the receptor backbone as rigid and approximates desolvation empirically. The present findings represent computational structural discoveries that require wet-lab verification prior to lead progression, including LanthaScreen/TR-FRET biochemical binding selectivity assays (TYK2 JH2 vs.\ JAK2 JH1), cellular inhibition of IFN$\alpha$-stimulated STAT1/STAT3 phosphorylation in human peripheral blood mononuclear cells (PBMCs), and counter-screening against EPO-stimulated JAK2/STAT5 signaling in TF-1 erythroleukemia cells.

\section{Methods}
\textbf{Receptor Preparation \& Alignment:} Crystal structures 6NZP (human TYK2 JH2, 2.35~Å) and 4IVA (human JAK2 JH1, 1.50~Å) were retrieved from RCSB PDB [2]. Waters and non-ligand crystallization additives were removed, chains were isolated to Chain~A, and polar hydrogens/protonation states were assigned at pH~7.4 via OpenBabel. Structures were aligned via Biopython \texttt{PairwiseAligner} and \texttt{Superimposer} across 18 conserved binding-site C$\alpha$ atoms ($\le 12.5$~Å\ from crystallographic ligand LB7), yielding a 1.887~Å\ C$\alpha$ RMSD.

\textbf{Ligand Triage \& Consensus Docking:} Compounds from ChEMBL~34 [3], control sets, and curated TYK2 analogs were filtered using RDKit \texttt{FilterCatalog} (PAINS, Brenk, NIH) [6], physicochemical bounds ($280 \le \text{MW} \le 490\text{ Da}$, $1.0 \le \text{cLogP} \le 4.8$, $45.0 \le \text{TPSA} \le 125.0\text{ Å}^2$, $\text{HBD} \le 4$, $\text{HBA} \le 9$, $\text{RotB} \le 8$), and SMARTS hinge-motif matching (\texttt{[n,N]\textasciitilde[c]\textasciitilde[n,N]}). Qualified molecules were embedded with ETKDGv3, energy-minimized with MMFF94s (\texttt{maxIters=400}), and converted to PDBQT via Meeko. Consensus docking was performed with AutoDock Vina 1.2.7 [4, 5] in a $22\times22\times22$~Å\ search box centered at $(14.16, -3.51, 28.73)$~Å\ (\texttt{exhaustiveness=16}, seeds 42, 137, 2026).

\textbf{Hardware Architecture:} Benchmarks were executed on an AWS EC2 \texttt{f2.6xlarge} instance equipped with a Xilinx UltraScale+ FPGA (\texttt{0000:34:00.0}, AFI \texttt{agfi-0b55f06c15a385187}) running Ubuntu 26.04 LTS. Direct userspace PCIe BAR0 MMIO access (\texttt{/sys/bus/pci/devices/0000:34:00.0/resource0}) was compiled via \texttt{gcc -O3 -march=native} (\texttt{liboikonomia\_voxel\_dma.so}) using \texttt{CLOCK\_MONOTONIC\_RAW} nanosecond timers and isolated CPU pinning (\texttt{taskset -c 2}).

\section{Data \& Code Availability}
All Python and C source code (\texttt{run\_tyk2\_jh2\_biorxiv.py}, \texttt{oikonomia\_voxel\_dma.c}), prepared PDBQT structures, PyMOL-ready 3D protein-ligand complexes, audited CSV hit tables, and hardware telemetry logs (\texttt{hardware\_telemetry.log}) are openly archived at \url{https://github.com/joeiko3424-svg/oikonomia-tyk2-jh2-fpga}.

\begin{thebibliography}{6}
\small
\bibitem{wrobleski2019} Wrobleski, S. T., Moslin, R., Lin, S., Zhang, Y., Spergel, S., Kempson, J., \dots \& Weinstein, D. S. (2019). Highly selective inhibition of tyrosine kinase 2 (TYK2) for the treatment of autoimmune diseases: discovery of the allosteric inhibitor BMS-986165. \textit{Journal of Medicinal Chemistry}, 62(20), 8973--8995.
\bibitem{berman2000} Berman, H. M., Westbrook, J., Feng, Z., Gilliland, G., Bhat, T. N., Weissig, H., \dots \& Bourne, P. E. (2000). The Protein Data Bank. \textit{Nucleic Acids Research}, 28(1), 235--242.
\bibitem{zdrazil2024} Zdrazil, B., Felix, E., Hunter, F., Manners, E. J., Blackshaw, J., Corbett, S., \dots \& Leach, A. R. (2024). The ChEMBL Database in 2023: a drug discovery platform spanning multiple bioactivity data types and time periods. \textit{Nucleic Acids Research}, 52(D1), D1180--D1192.
\bibitem{trott2010} Trott, O., \& Olson, A. J. (2010). AutoDock Vina: improving the speed and accuracy of docking with a new scoring function, efficient optimization, and multithreading. \textit{Journal of Computational Chemistry}, 31(2), 455--461.
\bibitem{eberhardt2021} Eberhardt, J., Santos-Martins, D., Tillack, A. F., \& Forli, S. (2021). AutoDock Vina 1.2.0: New docking methods, expanded force field, and Python bindings. \textit{Journal of Chemical Information and Modeling}, 61(8), 3891--3898.
\bibitem{landrum2024} Landrum, G. (2024). RDKit: Open-source cheminformatics. \url{https://www.rdkit.org}.
\end{thebibliography}

\end{document}